\documentclass[a4paper,10pt, twoside]{report} % Format A4, Schriftgröße 10, zweiseitige Ausrichtung
\usepackage[utf8]{inputenc}
\usepackage[ngerman]{babel} % Sprache: Deutsch nach neuer Rechtschreibung
\usepackage[T1]{fontenc}
\usepackage{lmodern} % andere Schriftart
\usepackage{geometry}
\usepackage{float}
\usepackage{amsmath,amsthm,amsfonts,amssymb} % Mathematik
\usepackage{color,graphicx,overpic}
\usepackage{listings} % Codefragmente
\usepackage{fancyhdr} % Kopf- und Fußzeilen
\usepackage[dvipsnames]{xcolor}
\usepackage{booktabs}
\usepackage[hidelinks,pdfencoding=auto]{hyperref} % Links & Verlinkung
\usepackage{longtable}
\usepackage{cite}
\usepackage[section]{placeins} % Bilder auf der richtigen Seite darstellen
\usepackage{enumitem}
\usepackage{subfiles}
\usepackage{parskip}
\usepackage{pgfplots}
\pgfplotsset{compat=1.8}
\setlength{\headheight}{28.7pt}
\setcounter{secnumdepth}{3}

\colorlet{mygray}{black!10}
\definecolor{ghostwhite}{rgb}{0.98, 0.98, 0.98}
\definecolor{bittersweet}{rgb}{0.96, 0.3, 0.26}
\definecolor{mygreen}{rgb}{0.0, 0.42, 0.24}

\definecolor{coralred}{rgb}{1.0, 0.25, 0.25}
\definecolor{deepfuchsia}{rgb}{0.76, 0.33, 0.76}
\definecolor{deeplilac}{rgb}{0.6, 0.33, 0.73}
\definecolor{glaucous}{rgb}{0.38, 0.51, 0.71}
\definecolor{pistachio}{rgb}{0.58, 0.77, 0.45}
\definecolor{green(ncs)}{rgb}{0.0, 0.62, 0.42}
\definecolor{mikadoyellow}{rgb}{0.93, 0.53, 0.18}
\definecolor{peridot}{rgb}{1.0, 0.75, 0.0}

\lstset{
    language=C++,
    backgroundcolor=\color{ghostwhite}, 
    basicstyle = \ttfamily,
    columns=fullflexible,
    breakatwhitespace=false,      
    breaklines=true,                
    captionpos=b,                    
    commentstyle=\color{mygreen}, 
    extendedchars=true,              
    frame=single,                   
    keepspaces=true,             
    keywordstyle=\color{blue},      
    language=c++,                 
    numbers=none,                
    numbersep=5pt,                   
    numberstyle=\tiny\color{blue}, 
    rulecolor=\color{mygray},        
    showspaces=false,               
    showtabs=false,                 
    stepnumber=1,                  
    stringstyle=\color{bittersweet},    
    tabsize=3,                      
    title=\lstname,
    numbers=left,
    columns = fullflexible,
    showstringspaces=false
}


\renewcommand{\lstlistingname}{Codeausschnitt}
\renewcommand{\lstlistlistingname}{Verzeichnis der Codeausschnitte}

\pdfinfo{
    /Title (SWP2021 AntiDoS - Finales Reviewdokument)
    /Creator (TeX)
    /Producer (pdfTeX 3.14)
    /Author (Fabienne Göpfert, Leon Leisten, Felix Hußlein, Robert Jeutter, Jakob Lerch, Tim Häußler, Johannes Lang, Tobias Scholz)
    /Subject (Abwehr von Denial-of-Service-Angriffen durch effiziente User-Space Paketverarbeitung)
}

%% Kopf- und Fußzeilen
\pagestyle{fancy}
\fancyhf{}
\fancyhead[EL]{\leftmark}
\fancyhead[OL]{\includegraphics[height=25pt]{img/aegis_logo_with_name.pdf}}
\fancyhead[OR]{\rightmark}
\fancyfoot[EL]{\thepage}
\fancyfoot[ER]{SWP2021 AntiDoS}
\fancyfoot[OL]{Finales Reviewdokument}
\fancyfoot[OR]{\thepage}
\renewcommand{\headrulewidth}{0.2pt} %fügt eine Linie unter der Kopfzeile ein

\renewcommand{\arraystretch}{1.3} %sorgt für größeren Abstand zwischen Zeilen in Tabellen

\begin{document}

\begin{titlepage} %Beginn der Titelseite
    \begin{center} %zentriert
        Technische Universität Ilmenau\\
        Fakultät Informatik und Automatisierung \\
        Fachgebiet Telematik/Rechnernetze\\
        %Technische Universität Ilmenau\\
        \vspace{2em}
        \includegraphics[width=0.35\linewidth]{img/tu_ilmenau_logo.png}\\
        %Technische Universität Ilmenau\\
        \vspace{2em}
        Finales Review zum Thema\\
        \vspace{1em}
        \Large\textsc{Abwehr von Denial-of-Service-Angriffen \\durch effiziente User-Space Paketverarbeitung}\\
        \vspace{0.5em}
        \Large\textsc{Codename: AEGIS}\\
        \vspace{0.5cm}
        \includegraphics[width=0.35\linewidth]{img/aegis_logo.pdf}\\
        \normalsize
        \vspace{2em}
        Autoren:\\
        \begin{table}[h]
            \centering
            \begin{tabular}{l r}
                Fabienne Göpfert & Tim Häußler    \\
                Felix Hußlein    & Robert Jeutter \\
                Johannes Lang    & Leon Leisten   \\
                Jakob Lerch      & Tobias Scholz  \\
            \end{tabular}
        \end{table}
        \vspace{1em}
        \normalsize {Betreuer:} Martin Backhaus\\
        \vspace{1em}
        Entwurf 21. Juli 2021
    \end{center}
\end{titlepage}
\pagebreak
\tableofcontents {\protect\thispagestyle{fancy}


\subfile{chapters/1-einleitung}
\subfile{chapters/2-grobentwurf}
\subfile{chapters/3-feinentwurf}
\subfile{chapters/4-tests}
\subfile{chapters/5-anforderungen}
\subfile{chapters/6-bugreview}
\subfile{chapters/7-statistiken}
\subfile{chapters/8-kimai}
\subfile{chapters/9-bewertung}
\subfile{chapters/10-abkuerzungsverzeichnis}
\subfile{chapters/11-glossar}


\bibliography{references}
\bibliographystyle{ieeetr}

\listoffigures {\protect\thispagestyle{fancy} % Abbildungsverzeichnis

\lstlistoflistings {\protect\thispagestyle{fancy}
\end{document}
